% =============================================================================
% Bloque 4: Estado del Arte
% Intro SMOTE/ADASYN + Intro LLMs + Dilema + Timeline + Insights + Posicionamiento
% =============================================================================

\section{Estado del arte}

% --- Slide: Intro SMOTE/ADASYN ---
\begin{frame}{Familia 1 — SMOTE y ADASYN: interpolación geométrica}
\begin{columns}[T]

\begin{column}{0.55\textwidth}
\small
\textbf{SMOTE} \cita{Chawla et al., 2002} y \textbf{ADASYN} \cita{He et al., 2008} son los métodos canónicos de oversampling.

\vspace{0.4em}
\textbf{Idea:}
\begin{itemize}\setlength\itemsep{2pt}
  \item Para cada muestra minoritaria $x_i$, encontrar sus $k$ vecinos más cercanos de la misma clase.
  \item Generar muestras nuevas sobre los segmentos:
\end{itemize}

\vspace{0.2em}
\centering
$x_{\text{new}} = x_i + \lambda \,(x_j - x_i), \quad \lambda \in [0,1]$
\raggedright

\vspace{0.5em}
\textbf{Diferencia clave:}
\begin{itemize}\setlength\itemsep{2pt}
  \item ADASYN \alert{pesa más} las muestras difíciles (rodeadas de otra clase).
\end{itemize}

\vspace{0.4em}
\good\ Geometría controlada. \quad \good\ Garantías formales.\\
\bad\ Genera \alert{vectores}, no texto. \quad \bad\ Solo rellena lo conocido.
\end{column}

\begin{column}{0.45\textwidth}
\centering
\vspace{0.3em}
\begin{tikzpicture}[scale=1.0]
  % Casco convexo
  \fill[methodSMOTE!10] (-1.3,-1.1) -- (1.3,-0.9) -- (1.5,0.7) -- (-0.5,1.3) -- (-1.3,-1.1);

  % Puntos reales
  \foreach \pt in {(-1.3,-1.1),(1.3,-0.9),(1.5,0.7),(-0.5,1.3),(0,0.2)}{\fill[methodSMOTE] \pt circle (3.5pt);}

  % Segmentos de interpolación
  \draw[methodSMOTE!60, dashed] (-1.3,-1.1) -- (1.3,-0.9);
  \draw[methodSMOTE!60, dashed] (-1.3,-1.1) -- (0,0.2);
  \draw[methodSMOTE!60, dashed] (0,0.2) -- (1.5,0.7);
  \draw[methodSMOTE!60, dashed] (-0.5,1.3) -- (0,0.2);

  % Sintéticos (en los segmentos)
  \foreach \pt in {(0,-1),(-0.65,-0.45),(0.75,0.45),(-0.25,0.75)}{\draw[uaiorange, very thick] \pt circle (3pt);}

  \node[font=\tiny, color=uaigray] at (0,-1.8) {{\color{methodSMOTE}\textbullet}\,reales \quad {\color{uaiorange}$\circ$}\,sintéticos};
  \node[font=\tiny\itshape] at (0,-2.1) {Interpolación entre $k$-NN};
\end{tikzpicture}
\end{column}

\end{columns}

\mensaje{Geometría como control: pero la salida es un vector, no una oración.}
\end{frame}

% --- Slide: Familia 2 — Aumentación textual ligera (entre SMOTE y LLMs) ---
\begin{frame}{Familia 2 — Aumentación textual ligera}
\small
Entre la \alert{interpolación geométrica} y la \alert{generación libre} hay un puente histórico: métodos que producen texto con modelos pequeños y operaciones controladas.

\vspace{0.5em}
\begin{columns}[T]

\begin{column}{0.33\textwidth}
\setbeamercolor{block title}{fg=white,bg=uaiblue}
\begin{block}{EDA}
\scriptsize
\cita{Wei \& Zou, 2019}\\[3pt]
4 ops a nivel \textbf{palabra}:\\
sinónimo, inserción, swap, deleción.\\[3pt]
\good\ Muy barato.\\
\bad\ Ruido en oraciones cortas.
\end{block}
\end{column}

\begin{column}{0.33\textwidth}
\setbeamercolor{block title}{fg=white,bg=uaiblue}
\begin{block}{Traducción inversa}
\scriptsize
\cita{Sennrich 2016; Edunov 2018}\\[3pt]
EN $\to$ otro idioma $\to$ EN.\\
Paráfrasis a nivel \textbf{oración}.\\[3pt]
\good\ Cambia sintaxis.\\
\bad\ Pierde matices.
\end{block}
\end{column}

\begin{column}{0.33\textwidth}
\setbeamercolor{block title}{fg=white,bg=uaiblue}
\begin{block}{BERT contextual}
\scriptsize
\cita{Kobayashi 2018; Wu 2019}\\[3pt]
Enmascarar tokens; BERT predice alternativas plausibles.\\[3pt]
\good\ Contexto semántico.\\
\bad\ Solo sustitución 1:1.
\end{block}
\end{column}

\end{columns}

\vspace{0.7em}
\begin{insight}
\scriptsize
Todos producen \textbf{texto}, pero con un \textbf{techo de creatividad bajo}: paráfrasis cercanas al original. Los LLMs rompen ese techo --- y por eso necesitan otro tipo de control.
\end{insight}
\end{frame}

% --- Slide: Familia 3 — LLMs ---
\begin{frame}{Familia 3 — LLMs: generación de texto bajo demanda}
\begin{columns}[T]

\begin{column}{0.58\textwidth}
\small
Pedirle a un LLM (GPT, Gemini, Claude\dots) que genere ejemplos nuevos para cada clase.

\vspace{0.3em}
\textbf{Receta típica} \cita{Whitehouse et al., 2023}:
\begin{enumerate}\setlength\itemsep{1pt}\scriptsize
  \item Prompt con descripción de clase y ejemplos reales.
  \item Pedir $N$ generaciones nuevas.
  \item Añadirlas al entrenamiento.
\end{enumerate}

\vspace{0.3em}
\textbf{Trabajos representativos:}
\begin{itemize}\setlength\itemsep{1pt}\scriptsize
  \item Suhaeni \& Yong (2024): $+12.7\%$ acc.\ en minoritarias.
  \item ImbLLM (2024): control de desbalance vía prompting.
  \item Cegin et al.\ (2025): diversidad guiada por LLM.
\end{itemize}

\vspace{0.4em}
{\scriptsize
\good\ Texto coherente. \quad \good\ Diversidad lingüística.\\
\bad\ \alert{Sin garantías geométricas}. \quad \bad\ Ruido y solapamiento.}
\end{column}

\begin{column}{0.42\textwidth}
\centering
\vspace{0.2em}
\begin{tikzpicture}[every node/.style={font=\scriptsize, align=center}]
  \node[draw=uaibluedark, fill=uaibluelight!15, rounded corners, very thick,
        text width=3.6cm, minimum height=1.0cm] (prompt) at (0, 2.7)
        {\scriptsize\textbf{Prompt}\\[1pt]
         \tiny\itshape ``Genera 50 mensajes de alegría\dots''};
  \draw[->, very thick, uaibluedark] (0, 2.05) -- (0, 1.65);
  \node[draw=methodOurs, fill=methodOurs!15, rounded corners, very thick,
        text width=2.4cm, minimum height=0.7cm] (llm) at (0, 1.2)
        {\textbf{LLM}};
  \draw[->, very thick, methodOurs] (0, 0.8) -- (0, 0.4);
  \node[draw=uaiorange, fill=uaiorange!10, rounded corners, very thick,
        text width=3.6cm, minimum height=1.0cm] (out) at (0, -0.2)
        {\scriptsize\textbf{$N$ oraciones nuevas}\\[1pt]
         \tiny algunas buenas, otras ruido\dots};
\end{tikzpicture}
\end{column}

\end{columns}

\mensaje{Texto sin control: los vectores caen donde quieran.}
\end{frame}

% --- Slide: El dilema (movido desde Introducción) ---
\begin{frame}{El dilema: SMOTE vs.\ LLMs lado a lado}
\begin{columns}[T]

\begin{column}{0.5\textwidth}
\centering
\setbeamercolor{block title}{fg=white,bg=methodSMOTE}
\begin{block}{SMOTE / interpolación}
\small
\good\ Geometría controlada\\
\good\ Garantías formales\\[2pt]
\bad\ Vectores sin texto coherente\\
\bad\ Solo rellena lo conocido
\end{block}
\vspace{0.4em}
\begin{tikzpicture}[scale=0.85]
  \fill[methodSMOTE!20] (0,0) circle (1.4);
  \foreach \pt in {(0,0.4),(0.7,0.2),(-0.5,0.6),(0.3,-0.5),(-0.7,-0.3),(0.6,-0.7)} {\fill[methodSMOTE] \pt circle (2pt);}
  \foreach \pt in {(0.4,0),(-0.2,0.1),(0.1,0.5),(-0.4,-0.5),(0.45,-0.2)} {\draw[uaiorange,thick] \pt circle (2pt);}
  \node[font=\tiny] at (0,-1.7) {Interpolación lineal};
\end{tikzpicture}
\end{column}

\begin{column}{0.5\textwidth}
\centering
\setbeamercolor{block title}{fg=white,bg=methodOurs}
\begin{block}{LLM (Gemini, GPT, Claude\dots)}
\small
\good\ Texto fluido y coherente\\
\good\ Diversidad lingüística\\[2pt]
\bad\ Geometría descontrolada\\
\bad\ Ruido y solapamiento entre clases
\end{block}
\vspace{0.4em}
\begin{tikzpicture}[scale=0.85]
  \fill[methodOurs!15] (0,0) circle (1.6);
  \foreach \pt in {(0,0.4),(0.7,0.2),(-0.5,0.6),(0.3,-0.5)} {\fill[methodOurs] \pt circle (2pt);}
  \foreach \pt in {(0.4,0),(-0.2,0.1),(0.1,0.5),(1.4,0.3),(-1.2,-0.4),(0.9,-1.0),(-0.8,1.1)} {\draw[uaiorange,thick] \pt circle (2pt);}
  \node[font=\tiny] at (0,-1.9) {Generación abierta};
\end{tikzpicture}
\end{column}

\end{columns}

\mensaje{¿Podemos quedarnos con lo mejor de ambos mundos?}
\end{frame}

% --- Slide: Timeline de 4 paradigmas ---
\begin{frame}{Cuatro paradigmas, una década de evolución}
\centering
\begin{tikzpicture}[every node/.style={font=\scriptsize}]
  % Línea de tiempo
  \draw[->, very thick, uaiblue!70] (0,0) -- (12,0) node[right, font=\scriptsize]{año};
  \foreach \x/\year in {0.5/2002, 3.5/2019, 6.5/2023, 9.5/2025}{
    \draw[uaiblue] (\x,-0.1) -- (\x,0.1);
    \node[below, font=\tiny] at (\x,-0.15) {\year};
  }

  % Paradigma 1
  \node[draw=uaiblue, fill=uaibluelight!15, rounded corners, text width=2.4cm, align=center,
        minimum height=1.2cm, font=\tiny] at (1.5, 1.5) (p1)
        {\textbf{Geométrico clásico}\\SMOTE, ADASYN, EDA};
  \draw[->, uaiblue!50] (0.5,0) -- (p1);

  % Paradigma 2
  \node[draw=uaiblue, fill=uaibluelight!15, rounded corners, text width=2.4cm, align=center,
        minimum height=1.2cm, font=\tiny] at (4.5, 1.5) (p2)
        {\textbf{Embeddings profundos}\\mixup, MixText, SMOTExT};
  \draw[->, uaiblue!50] (3.5,0) -- (p2);

  % Paradigma 3
  \node[draw=uaiblue, fill=uaibluelight!15, rounded corners, text width=2.4cm, align=center,
        minimum height=1.2cm, font=\tiny] at (7.5, 1.5) (p3)
        {\textbf{Generación LLM}\\GPT-3/4, Gemini, Claude};
  \draw[->, uaiblue!50] (6.5,0) -- (p3);

  % Paradigma 4
  \node[draw=uaiorange, very thick, fill=uaiorange!15, rounded corners, text width=2.4cm, align=center,
        minimum height=1.2cm, font=\tiny] at (10.5, 1.5) (p4)
        {\textbf{Híbridos + filtros}\\\textit{esta tesis}};
  \draw[->, uaiorange] (9.5,0) -- (p4);

  % Notas debajo
  \node[font=\tiny, text width=2.6cm, align=center] at (1.5,-1) {Chawla 2002\\Wei \& Zou 2019};
  \node[font=\tiny, text width=2.6cm, align=center] at (4.5,-1) {Zhang 2018\\Bayer 2023};
  \node[font=\tiny, text width=2.6cm, align=center] at (7.5,-1) {Suhaeni \& Yong\\\textbf{+12.7\%} acc};
  \node[font=\tiny, text width=2.6cm, align=center] at (10.5,-1) {Cegin 2025, ImbLLM\\Springer 2025};
\end{tikzpicture}

\vfill
\mensaje{Cada paradigma resuelve un problema y crea otro nuevo.}
\end{frame}

% --- Slide: Tres deficiencias (después del Timeline) ---
\begin{frame}{Tres deficiencias en los métodos actuales}
\small
\begin{columns}[T]

\begin{column}{0.33\textwidth}
\setbeamercolor{block title}{fg=white,bg=uaiorange}
\begin{block}{1. Sin validación}
Las muestras del LLM se aceptan \textbf{todas}, sin verificar si caen en la región correcta del espacio.
\vspace{0.3em}\\
{\scriptsize Texto gramatical $\neq$ texto en la zona correcta.}
\end{block}
\end{column}

\begin{column}{0.33\textwidth}
\setbeamercolor{block title}{fg=white,bg=uaiamber}
\begin{block}{2. Peso uniforme}
Todas las muestras pesan igual en el clasificador, ignorando la \textbf{calidad geométrica} de cada una.
\vspace{0.3em}\\
{\scriptsize Una en zona densa = una en zona dispersa.}
\end{block}
\end{column}

\begin{column}{0.33\textwidth}
\setbeamercolor{block title}{fg=white,bg=uaiblue}
\begin{block}{3. Un solo dataset}
Casi todos los estudios validan en \textbf{un único} corpus o tarea.
\vspace{0.3em}\\
{\scriptsize ¿Generaliza el filtrado entre dominios y tareas?}
\end{block}
\end{column}

\end{columns}

\vfill
\begin{tikzpicture}[overlay,remember picture]
\end{tikzpicture}

\mensaje{La señal geométrica está disponible y nadie la está usando.}
\end{frame}

% --- Slide: Gap metodológico + idea conceptual ---
\begin{frame}{El gap metodológico: filtrar la salida del LLM}
\begin{columns}[T]

\begin{column}{0.52\textwidth}
\small
Las deficiencias anteriores apuntan a una intervención conceptualmente simple, no estudiada sistemáticamente.

\vspace{0.4em}
\textbf{La idea:}
\begin{enumerate}\setlength\itemsep{1pt}\scriptsize
  \item Generar $N$ muestras con el LLM (como hoy).
  \item Proyectarlas al espacio de embeddings.
  \item Medir qué tan cerca caen del \textbf{cluster real} de su clase.
  \item \alert{Rechazar las que se alejan demasiado}.
\end{enumerate}

\vspace{0.4em}
\textbf{¿Por qué tiene sentido?}
\begin{itemize}\setlength\itemsep{1pt}\scriptsize
  \item La geometría \alert{ya está ahí}: usamos embeddings para el clasificador.
  \item Cero costo extra de generación; el filtro opera \emph{a posteriori}.
  \item Agnóstico al LLM concreto.
\end{itemize}
\end{column}

\begin{column}{0.48\textwidth}
\centering
\vspace{0.2em}
\begin{tikzpicture}[scale=0.85]
  % Cluster real (fondo)
  \fill[methodSMOTE!18] (0,0) circle (1.5);

  % Frontera del filtro (dashed)
  \draw[uaibluedark, dashed, thick] (0,0) circle (1.9);

  % Puntos reales (sólidos)
  \foreach \pt in {(0,0.2),(0.4,0.4),(-0.5,0.3),(0.6,-0.4),(-0.3,-0.5),(0.2,0.7),(-0.6,-0.2),(0.7,0.1)} {\fill[methodSMOTE] \pt circle (3pt);}

  % LLM aceptados (verdes, dentro del filtro)
  \foreach \pt in {(0.9,0.5),(-0.8,-0.7),(0.3,1.0),(-0.4,1.1),(1.0,-0.5),(-1.1,0.4)} {\draw[methodOurs, very thick] \pt circle (3pt);}

  % LLM rechazados (rojos, fuera del filtro)
  \foreach \pt in {(2.4,0.7),(-2.3,0.9),(2.0,-1.4),(-1.7,-2.0),(0.1,2.3),(2.5,-0.3)} {\draw[uaiorange, very thick] \pt circle (3pt);}

  % Leyenda compacta
  \node[font=\tiny, color=methodSMOTE, anchor=west] at (-2.8,3.0) {{\fontsize{6}{6}\selectfont$\bullet$}\,reales};
  \node[font=\tiny, color=methodOurs, anchor=west] at (-2.8,2.7) {$\circ$\,aceptado};
  \node[font=\tiny, color=uaiorange, anchor=west] at (-2.8,2.4) {$\circ$\,rechazado};
  \node[font=\tiny, color=uaibluedark, anchor=west] at (-2.8,2.1) {- -\,filtro};
\end{tikzpicture}
\end{column}

\end{columns}

\mensaje{La señal geométrica ya existe — falta usarla para decidir qué entra al entrenamiento.}
\end{frame}

% --- Slide: Posicionamiento ---
\begin{frame}{Dónde encaja esta propuesta}
\begin{center}
\begin{tikzpicture}[scale=0.72, every node/.style={font=\scriptsize}, transform shape]
  \fill[uaiblue!25, opacity=0.7] (-1.4,0) circle (1.9);
  \fill[uaigreen!25, opacity=0.7] (1.4,0) circle (1.9);
  \fill[uaiamber!25, opacity=0.7] (0,-1.9) circle (1.9);

  \node[font=\bfseries\small, color=uaiblue, text width=2.4cm, align=center] at (-3.2,1.4) {Validación post-generación};
  \node[font=\bfseries\small, color=uaigreen!70!black, text width=2.4cm, align=center] at (3.2,1.4) {Ponderación adaptativa};
  \node[font=\bfseries\small, color=uaiamber!70!black, text width=2.6cm, align=center] at (0,-3.4) {Validación multi-tarea (NER)};

  \node[font=\tiny, color=uaiblue!70!black] at (-2.2,0) {LOF};
  \node[font=\tiny, color=uaigreen!60!black, text width=1.9cm, align=center] at (2.2,0) {Importance Sampling};
  \node[font=\tiny, color=uaiamber!70!black] at (0,-1.3) {MELM 2022};

  \node[draw=uaiorange, very thick, fill=white, rounded corners,
        font=\bfseries\scriptsize, text width=1.9cm, align=center,
        minimum height=0.8cm] at (0,-0.3) {Esta propuesta};
\end{tikzpicture}
\end{center}

\mensaje{Contribución: integrar las tres dimensiones por primera vez.}
\end{frame}
